% Original finite bridge prepared for the Robin/iPRA report.
% This is a fragment intended for inclusion in the parent's document.
\section{An explicit finite least-common-multiple bridge}

This section proves an elementary reduction which does not use the Riemann
hypothesis or an oscillation theorem. Its input is one strict inequality at
a specified finite frontier. Its output is one explicitly constructed integer
violating Robin's inequality. All arithmetic sums and products are bounded.

For an integer $P\ge2$, define
\[
 \vartheta(P)=\sum_{\substack{p\le P\\p\text{ prime}}}\log p,
 \qquad
 a_p(P)=\max\{a\le P:p^a\le P\},
 \qquad
 \psi(P)=\sum_{\substack{p\le P\\p\text{ prime}}}a_p(P)\log p.
\]
The maximum defining $a_p(P)$ is a bounded search. Put
\[
 L_P=\operatorname{lcm}(1,\ldots,P),\qquad
 A(n)=\frac{\sigma(n)}n,
 \qquad
 M(n)=\log A(n)-\gamma-\log\log\log n.
\]
The last expression is used only for $n>e$. Since the logarithm is strictly
increasing, $M(n)>0$ is equivalent to
$\sigma(n)>e^\gamma n\log\log n$.

For $\vartheta(P)>1$, define
\[
 F(P)=e^\gamma\log\vartheta(P)
       \prod_{\substack{p\le P\\p\text{ prime}}}\left(1-\frac1p\right),
\]
and the two nonnegative losses
\[
 D(P)=\log\log\psi(P)-\log\log\vartheta(P),
 \qquad
 T(P)=\sum_{\substack{p\le P\\p\text{ prime}}}
       -\log\left(1-p^{-a_p(P)-1}\right).
\]

\begin{lemma}[Exact finite identity]
If $\vartheta(P)>1$, then
\[
 M(L_P)=-\log F(P)-D(P)-T(P).
\]
\end{lemma}
\begin{proof}
Unique prime factorization gives
$L_P=\prod_{p\le P}p^{a_p(P)}$, and hence $\log L_P=\psi(P)$.
The finite geometric-series formula and multiplicativity of $\sigma$ give
\[
 A(L_P)=\prod_{p\le P}
        \frac{1-p^{-a_p(P)-1}}{1-p^{-1}}.
\]
Taking logarithms yields
$\log A(L_P)=-T(P)-\sum_{p\le P}\log(1-p^{-1})$.
On the other hand,
\[
 \log F(P)=\gamma+\log\log\vartheta(P)
              +\sum_{p\le P}\log(1-p^{-1}).
\]
Substituting these equalities into the definition of $M(L_P)$ proves the
claim. Every logarithm is in its positive domain: $\psi(P)\ge\vartheta(P)>1$.
\end{proof}

\begin{lemma}[Losses at a square frontier]
Let $m\ge4$ be an integer, put $P=m^2$, and suppose
$\vartheta(P)\ge P/4$. Then
\[
 0\le D(P)\le\frac8m,\qquad
 0\le T(P)\le\frac4m,
 \qquad
 D(P)+T(P)\le\frac{12}m.
\]
\end{lemma}
\begin{proof}
Only primes $p\le m$ contribute to $\psi(P)-\vartheta(P)$.
For each such prime, $(a_p(P)-1)\log p\le\log P$; there are at most
$m$ such primes. Therefore
\[
 0\le\psi(P)-\vartheta(P)\le m\log P.
\]
For $v\ge u>1$, two applications of $\log(1+t)\le t$ give
\[
 \log\log v-\log\log u
 \le\frac{\log v-\log u}{\log u}
 \le\frac{v-u}{u\log u}.
\]
Since $P\ge16$,
$\log(P/4)\ge\tfrac12\log P$. Consequently
\[
 D(P)\le\frac{m\log P}{(P/4)\log(P/4)}\le\frac8m.
\]

For $0\le u\le1/2$, the geometric-series estimate gives
$-\log(1-u)\le u/(1-u)\le2u$.
If $p\le m$, maximality of $a_p(P)$ implies $p^{a_p(P)+1}>P$, so
$p^{-a_p(P)-1}<1/P$. These at most $m$ summands contribute at most $2/m$.
If $m<p\le P$, then $a_p(P)=1$, so the remaining summands contribute at most
\[
 2\sum_{j=m+1}^{P}\frac1{j^2}
 \le2\sum_{j=m+1}^{P}\frac1{j(j-1)}
 =2\left(\frac1m-\frac1P\right)\le\frac2m.
\]
Here the inequality is termwise and the equality is a finite telescoping
sum. This proves $T(P)\le4/m$, and addition proves the last claim.
\end{proof}

\begin{lemma}[A rough elementary Chebyshev lower bound]
For every integer $m\ge16$,
\[
 \vartheta(m^2)\ge\frac{m^2}4.
\]
\end{lemma}
\begin{proof}
Put $P=m^2$ and $h=\lfloor P/2\rfloor$.
Legendre's formula shows that the exponent of a prime $p$ in
$\binom{2h}{h}$ is
\[
 \sum_{j\ge1}\left(\left\lfloor\frac{2h}{p^j}\right\rfloor
       -2\left\lfloor\frac h{p^j}\right\rfloor\right).
\]
The sum is bounded by $j\le2h$. Every summand is either zero or one,
and it vanishes for $p^j>2h$. Thus the exponent is at most
$a_p(2h)$, proving $\binom{2h}{h}\mid L_{2h}$.
The central coefficient is the largest of the $2h+1$ coefficients in
$(1+1)^{2h}$; hence
\[
 \psi(P)\ge\log\binom{2h}{h}
 \ge2h\log2-\log(2h+1)
 \ge(P-1)\log2-\log(P+1).
\]

The logarithm series at $1/3$ gives $2/3\le\log2\le3/4$.
Also, since $m/16\ge1$ and $\log t\le t-1$,
\[
 \log m=4\log2+\log(m/16)
 \le3+(m/16-1)=2+m/16\le3m/16.
\]
Therefore $\log P=2\log m\le3m/8$, while
\[
 \log(P+1)\le\log(2P)\le\frac34+\frac{3m}8.
\]
The previous prime-power count gives
$\psi(P)-\vartheta(P)\le m\log P\le3m^2/8$. Combining the estimates,
\[
 \vartheta(m^2)
 \ge\frac{7m^2}{24}-\frac{3m}8-\frac{17}{12}
 =\frac{m^2}4+\frac{m^2-9m-34}{24}.
\]
Finally $m^2-9m-34=(m-16)(m+7)+78\ge0$ for $m\ge16$.
\end{proof}

\begin{theorem}[Explicit finite counterexample transfer]
Let $m\ge16$ be an integer. If
\[
 -\log F(m^2)>\frac{12}m,
\]
then the explicitly computed integer $N=L_{m^2}$ satisfies $N>5040$ and
\[
 \sigma(N)>e^\gamma N\log\log N.
\]
\end{theorem}
\begin{proof}
The previous two lemmas and the exact identity give
\[
 M(N)=-\log F(m^2)-D(m^2)-T(m^2)>0.
\]
Moreover $m^2\ge16$ and $L_{16}=720720$ divides $L_{m^2}$, so $N>5040$.
The equivalence between $M(N)>0$ and the displayed strict inequality was
already established.
\end{proof}

\begin{lemma}[Passing to a square frontier]
Let $Q\ge256$ be an integer, $m=\lfloor\sqrt Q\rfloor$, and $P=m^2$.
Then
\[
 \log F(P)\le\log F(Q)+\frac4m.
\]
Consequently $-\log F(Q)>16/m$ yields a Robin counterexample at $L_P$.
\end{lemma}
\begin{proof}
We have $m\ge16$ and $0\le Q-P\le2m$.
The term $\log\log\vartheta$ is nondecreasing, so
\[
 \log F(P)-\log F(Q)
 \le\sum_{P<p\le Q}-\log(1-p^{-1})
 \le\frac{2(Q-P)}{P+1}\le\frac4m.
\]
The previous theorem now applies.
\end{proof}

This lemma makes the connection with the conventional nonsquare oscillation
statement explicit. Suppose $c>0$, $0<b<1/2$, and an integer $Q\ge256$
satisfies $-\log F(Q)\ge cQ^{-b}$. Set $\delta=1/2-b$.
Since $m\ge\sqrt Q/2$, the additional condition $cQ^\delta>32$ implies
$-\log F(Q)>16/m$ and hence a counterexample. For rational $c,b$, choose
integers $B\ge\max(2,64/c)$ and $k\ge1/\delta$. Then the sufficient
cutoff $Q\ge\max(256,B^k)$ is primitive recursive in their rational codes.
Classically, Nicolas's excursions occur beyond every cutoff, so this
finishes the transfer from that analytic input. It does not itself prove
the analytic input or its iPRA interpretation.

The result also holds for $m\ge4$ when $\vartheta(m^2)\ge m^2/4$ is supplied
as an explicit hypothesis. The choice $m\ge16$ removes that hypothesis by
the completely elementary proof above. A supplied rational positive
gap determines a primitive-recursive precision for the log enclosure
and therefore a finite certificate. Strictness by itself is not being
used as an unbounded-search termination bound. The proof contains no
search for the input frontier.

\paragraph{What this does and does not finish.}
An analytic oscillation theorem would still have to supply an appropriate
finite frontier. Classical Nicolas--Landau oscillation gives such excursions
arbitrarily far out, of size $cP^{-b}$ for some $0<b<1/2$. The finite loss
is only of order $P^{-1/2}$, so that implication follows classically. To
obtain an iPRA theorem, the analytic result needs an explicitly justified
interpretation using bounded arithmetic computations and the permitted
induction. Replacing that obligation by the assertion that a classical Lean
proof exists would not prove it.
